\section{The remaining tail carries almost all actual boundary mass}
\label{fml-section}

Keep the actual roots $w_k=3k+\zeta$ for $B\le k<2B$, where
$B,n$ are positive integers. Put
\[
 Q_k=9k^2+3k+1,\qquad T_n(k)=|P(w_k^n)|,\qquad
 t_k=\log c_n(k).
\]
Their primitive unramified powers have nonzero boundary, and
\begin{equation}\label{fml-heights}
 t_k\ge n\log(3B),\qquad
 \log|P(w_k^d)|\le3d\log(6B+1).
\end{equation}
Here the observable is full valuation mass, including shallow primes.
It is different from both positive excess and signed cost.

\begin{theorem}[Full-mass mean with rank spacing]\label{fml-mean}
For real $Z\ge5$ with $Z>3n$, define
$H_Z(k)=\sum_{5\le q\le Z}v_q(T_n(k))\log q$, over primes.
Then
\begin{align}
 \frac1B\sum_{B\le k<2B}\frac{H_Z(k)}{t_k}
 \le{}&\frac{6\tau(n)\log Z(1+\log Z)}{n\log(3B)}\notag\\
      &+\frac{36Z\sigma(n)}{Bn\log(Z/(3n))}
          +\frac{\log n}{n\log(3B)}.
          \label{fml-mean-bound}
\end{align}
The first $\tau(n)$ can be replaced by $\sum_{d\mid n}\varphi(d)/d$.
\end{theorem}
\begin{proof}
Partition supported actual roots at $q\ge5$ by their rank
$d=\operatorname{ord}((w_k/\bar w_k)^3\bmod q)$. It divides $n$,
and its first depth $s=v_q(T_d(k))$ satisfies
$v_q(T_n(k))=s+v_q(n)$. Primes dividing $Q_k$ have no boundary
contribution. The total additional lifting is at most $\log n$
at each root.

For fixed $d,q$, the established exact-rank count has at most
$3\varphi(d)$ simple roots, each lifting uniquely. Thus the actual
interval count at depth at least $e$ is at most
$3\varphi(d)(B/q^e+1)$. Sum from $e=1$, with the cap
$\lfloor3d\log(6B+1)/\log q\rfloor$ from \eqref{fml-heights}.
The full first-depth contribution, summed over the actual block, is
at most
\begin{equation}\label{fml-prime-rank}
 \frac{3\varphi(d)B\log q}{q-1}
                      +9\varphi(d)d\log(6B+1).
\end{equation}
This remains valid if the cap is zero; no actual simplicity of depth
is assumed.

Eligible primes have $q=3dj+1$ or $q=3dj-1$, $j\ge1$.
In the first progression $1/(q-1)=1/(3dj)$; in the second it is
at most $1/(dj)$, since $3dj-2\ge dj$. All relevant $j$ are at
most $Z$, and the harmonic sum gives
\[
 \sum_{\substack{q\le Z\\q\text{ eligible at }d}}
        \frac{\log q}{q-1}
 \le\frac{2\log Z(1+\log Z)}d.
\]
The first part of \eqref{fml-prime-rank}, summed over $d\mid n$
and divided by $Bn\log(3B)$, gives the first term of
\eqref{fml-mean-bound}. We used $\varphi(d)/d\le1$.

For the endpoint part, the previously stated Brun--Titchmarsh bound
counts at most $4Z/[\varphi(3d)\log(Z/(3d))]$ eligible primes.
Its premise holds for every divisor since $Z>3n$.
Using $\varphi(3d)\ge2\varphi(d)$ gives a total endpoint cost at
most $18Z\sigma(n)\log(6B+1)/\log(Z/(3n))$.
Divide by $Bn\log(3B)$ and use $\log(6B+1)\le2\log(3B)$.
Finally append the lifting contribution. All normalized costs are
nonnegative, so the lower height in \eqref{fml-heights} is sufficient.
This proves the theorem using actual interval counts.
\end{proof}

\begin{theorem}[Actual full mass beyond the all-index window]
\label{fml-concentration}
Let $n\to\infty$ through all positive integers, and put
\[
 B=n^4,\qquad
 Z=\left\lfloor\frac{B\sqrt{\log n}}{1+\log\log n}\right\rfloor,
 \qquad \eta_n=(\log n)^{-1/4}.
\]
There are absolute constants $C,N$ such that for every $n\ge N$,
\begin{equation}\label{fml-all-prime-mean}
 \frac1B\sum_{B\le k<2B}
    \frac{\displaystyle\sum_{q\le Z}v_q(T_n(k))\log q}{t_k}
 \le C(\log n)^{-1/2}.
\end{equation}
This sum includes $q=2,3$. Outside a proportion at most $C\eta_n$
of the actual roots,
\begin{equation}\label{fml-tail-three}
 3-A_1\frac{\log(4n)}n-\eta_n
 \le\frac{\displaystyle\sum_{q>Z}v_q(T_n(k))\log q}{t_k}
 \le3.
\end{equation}
If $n$ tends to infinity through primes at least five, the mass
restricted further to exact rank $n$ obeys, on that same set,
\begin{equation}\label{fml-top-three}
 \frac{\displaystyle\sum_{\substack{q>Z\\d_q(k)=n}}
                    v_q(T_n(k))\log q}{t_k}\longrightarrow3.
\end{equation}
At these primes the valuation is exactly the first depth.
\end{theorem}
\begin{proof}
The elementary inequality $\tau(n)\le2\sqrt n$ makes the first
term in \eqref{fml-mean-bound} $O(\log n/\sqrt n)$.
Lemma~\ref{aw-sigma} bounds $\sigma(n)/n$ by
$e^{18}(1+\log\log n)$; its extra factor is canceled by the same
factor in the denominator of $Z/B$. Since
$\log(Z/(3n))\sim3\log n$, the endpoint term is
$O((\log n)^{-1/2})$, and the lifting term is $O(1/n)$.

The all-prime theorem with unit residual and $\lambda=1/n$ gives
\[
 \frac{v_2(T_n)\log2+v_3(T_n)\log3}{t_k}
 \le A_2(4\log2+9\log3)\frac{\log^2(4n)}n.
\]
Sector normalization meets its positive-coordinate domain and preserves
these valuations and the height. This proves
\eqref{fml-all-prime-mean}. Markov's inequality at $\eta_n$ leaves
at most the stated exceptional proportion. For every remaining root,
subtract its complete small-prime mass from
\[
 \frac{\log T_n}{t_k}=3-\frac{\Delta_n}{t_k},\qquad
 0\le\frac{\Delta_n}{t_k}\le A_1\frac{\log(4n)}n.
\]
This proves \eqref{fml-tail-three}.

When $n$ is prime, only ranks one and $n$ occur. Above $Z>3n$
there is no exponent lifting. The rank-one full mass is bounded by
$\log T_1(k)$, and
\[
 \frac{\log T_1(k)}{t_k}
 \le\frac{(3/2)\log Q_k}{(n/2)\log Q_k}=\frac3n.
\]
Removing it from \eqref{fml-tail-three} proves
\eqref{fml-top-three}, uniformly on the nonexceptional set.
\end{proof}

This locates the main amount of integer boundary size in the actual
remaining tail. At prime index, it is the first-depth top-rank packet,
not an old lower-rank contribution. It does not supply a lower bound
on the signed cost: a prime of depth one contributes positive full
mass but signed credit $-2\log q$. Almost all the mass above the
window could have that useful shallow form. Neither its radical nor
its deep private owners are controlled by the theorem.

The proof is an ordinary analytic argument using the established
actual rank/lifting counts, the Brun--Titchmarsh estimate already
used in Theorem~\ref{sw-estimate}, and the two-place theorem only
for the fixed primes two and three and the archimedean place.
It uses no modularity or probabilistic independence and is not a
claim of a Lean formalization. Exceptional roots and the net signed
tail remain open.
